\chapter*{Acknowledgements}
\renewcommand{\baselinestretch}{\linespacing}
\small\normalsize

I would like to thank my supervisor Christoph Flamm for the opportunity and his patience.
For tolerating my occasional slow grasp of concepts and for all the support I received, I am grateful.

Gratitude is also extended to the entire TBI working group for their support and for creating a stimulating environment. Special thanks to Manuel Uhlir for getting me started and for the many helpful discussions and suggestions along the way.

Heartfelt thanks are due to my mother and grandmother, whose constant faith in me has carried me through many challenges. I also wish to acknowledge my grandfather for the memorable time I was able to share with him.

Finally, thank you, Aila, for your warmth and kindness at every step of this journey.

\newpage
\pdfbookmark[0]{Contents}{contents_bookmark}
\renewcommand{\baselinestretch}{1.3}\normalsize
\tableofcontents
\renewcommand{\baselinestretch}{1.5}\normalsize
\clearpage
\chapter*{Abstract}
%min. 500 Zeichen inkl. Leerzeichen

Mass spectrometry is a cornerstone of analytical chemistry, metabolomics, and pharmaceutical research. Machine learning increasingly improves forward prediction of spectra from known structures, but the inverse problem --- inferring plausible molecular structures from a spectrum --- remains difficult, ambiguous, and hard to interpret.

This work introduces intermediate fragments generated by graph-transformation tools as an explicit fragmentation layer for bidirectional structure $\leftrightarrow$ spectrum translation, making the process more transparent than black-box models.
The main contribution is the framework itself, together with a mass-controlled evaluation that separates learned structural signal from mere mass consistency.
In a proof-of-concept study the forward model predicts molecule-specific spectra and spectrum-to-structure retrieval is assessed, giving a concrete basis for future work on learned fragment generation.

\phantomsection
\addcontentsline{toc}{chapter}{Abstract}
\clearpage
\chapter*{Zusammenfassung}
%min. 500 Zeichen inkl. Leerzeichen
Massenspektrometrie ist eine Schlüsseltechnik in der analytischen Chemie, der Metabolomik und der pharmazeutischen Forschung. Methoden des maschinellen Lernens verbessern zunehmend die Vorwärtsvorhersage von Spektren aus bekannten Strukturen; das inverse Problem --- die Ableitung plausibler Molekülstrukturen aus einem Spektrum --- bleibt jedoch schwierig, mehrdeutig und schwer interpretierbar.

Diese Arbeit führt Zwischenfragmente, die durch graphbasierte Transformationstools erzeugt werden, als explizite Fragmentierungsschicht für die bidirektionale Struktur-$\leftrightarrow$-Spektrum-Übersetzung ein und macht den Prozess dadurch transparenter als Black-Box-Modelle.
Der wesentliche Beitrag ist das Framework selbst sowie eine massenkontrollierte Evaluation, die gelerntes strukturelles Signal von bloßer Massenkonsistenz trennt.
In einer Machbarkeitsstudie (Proof of Concept) sagt das Vorwärtsmodell molekülspezifische Spektren vorher und die Spektrum-zu-Struktur-Wiederauffindung (Retrieval) wird bewertet; dies bildet eine konkrete Grundlage für zukünftige Arbeiten zur gelernten Fragmenterzeugung.

\phantomsection
\addcontentsline{toc}{chapter}{Zusammenfassung}
\listoftables
\phantomsection
\addcontentsline{toc}{chapter}{List of Tables}
\listoffigures
\phantomsection
\addcontentsline{toc}{chapter}{List of Figures}
